\begin{ejercicio}
    % ID: 77
Consideremos el \textit{complejo de cadenas aumentado} de $X$:
\[
\begin{tikzcd}[column sep=1.7em]
\cdots \arrow[r] & S_q(X) \arrow[r] & S_{q-1}(X) \arrow[r] & \cdots \arrow[r] & S_2(X) \arrow[r, "\partial_2"] & S_1(X) \arrow[r, "\partial_1"] & S_0(X) \arrow[r, "\varepsilon"] & R \arrow[r] & 0.
\end{tikzcd}
\]
Definiendo $\widetilde{H}_q(X; R)$ como la homología de este complejo de cadenas aumentado, se obtiene
\begin{align*}
\widetilde{H}_0(X; R) &\equiv \ker\varepsilon / \operatorname{Im}\partial_1, \\
\widetilde{H}_q(X; R) &= H_q(X; R) \qquad \text{si } q > 0.
\end{align*}
Ver que esta definición es isomorfa a cualquiera de las dos anteriores (recordemos que $\operatorname{Im}\partial_1 \subset \ker\varepsilon$).
\end{ejercicio}

\begin{demostracion}

\end{demostracion}
